\documentclass{ctexart}

% 支持从项目根目录或当前目录编译。
\makeatletter
\def\input@path{{../../}}
\makeatother
\usepackage{researchnotes}

\title{根式数列的收敛性}
\author{}
\date{}

\begin{document}
\maketitle

\section{题目与解法}

\begin{example}
  设实数列 $\{x_n\}_{n\ge1}$ 满足
  \[
    x_n=\sum_{k=1}^{n}\frac{1}{\sqrt{k}}-2\sqrt{n}.
  \]
  证明 $\displaystyle\lim_{n\to\infty}x_n$ 存在。
\end{example}

\paragraph{思路}
利用\keyterm{单调有界定理}（monotone convergence theorem）：
单调递减且有下界的实数列必收敛。
先将相邻两项的差有理化以判断单调性，再利用根式之差构造可以相消求和的下界。

\begin{proof}
\color{black}
首先，对任意整数 $n\ge1$，有
\begin{align*}
  x_{n+1}-x_n
  &=\frac{1}{\sqrt{n+1}}-2\bigl(\sqrt{n+1}-\sqrt{n}\bigr)\\
  &=\frac{1}{\sqrt{n+1}}-
    \frac{2}{\sqrt{n+1}+\sqrt{n}}<0.
\end{align*}
最后一个不等式成立，是因为
$\sqrt{n+1}+\sqrt{n}<2\sqrt{n+1}$，且各分母均为正数。
因此，$\{x_n\}$ 严格单调递减。

其次，对任意整数 $k\ge1$，由
$\sqrt{k+1}+\sqrt{k}\ge2\sqrt{k}>0$，得到
\[
  \frac{1}{\sqrt{k}}
  \ge\frac{2}{\sqrt{k+1}+\sqrt{k}}
  =2\bigl(\sqrt{k+1}-\sqrt{k}\bigr).
\]
两边从 $k=1$ 到 $n$ 求和，利用相邻项相消，得到
\[
  \sum_{k=1}^{n}\frac{1}{\sqrt{k}}
  \ge2\sum_{k=1}^{n}\bigl(\sqrt{k+1}-\sqrt{k}\bigr)
  =2\sqrt{n+1}-2.
\]
于是
\[
  x_n\ge2\bigl(\sqrt{n+1}-\sqrt{n}\bigr)-2>-2.
\]
因此，$\{x_n\}$ 有下界。

综上，$\{x_n\}$ 单调递减且有下界，由单调有界定理，
$\displaystyle\lim_{n\to\infty}x_n$ 存在且为有限实数。
\end{proof}

\end{document}
